\section{问题分析}

题目围绕温度响应、催化剂组合因素、收率优化及新增实验设计展开。本文已冻结的问题一只处理前两类实验记录中的温度响应和时间变化，不对组合间差异作归因或优劣比较；后续问题若需比较催化剂因素或优化条件，应另行建立模型，不能反向改变本问的组内经验关系。

\subsection{问题一分析}

附件1的数学对象是同一催化剂组合内部的稀疏温度--响应曲线，而非 114 条可混合的同质样本。每组仅有 5--7 个实际温度点，因此以线性模型给出最低复杂度的总体趋势，并以只多一个参数的二次多项式检验当前温区内是否存在可辨识曲率。两者均为经验近似，不代表真实反应机理；选用冻结的组内留一法比较预测误差，并同时检查实验温度范围内的物理边界。

附件2只有同一次实验的 7 个不等间隔时间观测，且未知其 A/B 催化剂编号。故直接保持时间顺序，计算转化率、选择性、派生收率和相邻区间变化率；线性趋势仅用于量化给定时间窗内的平均转化率变化，不作长期外推或失活机理判定。

\subsection{问题二分析}

\todo{待后续已确认成果写入。问题二不得以问题一的组内经验拟合替代因素影响分析。}

\subsection{问题三分析}

\todo{待后续已确认成果写入。问题三的优化结论不得从问题一的二次曲线顶点直接推得。}
